\documentclass[11pt]{article}

\usepackage[margin=1in]{geometry}
\usepackage{amsmath,amssymb,bm}
\usepackage{booktabs}
\usepackage{siunitx}
\usepackage{hyperref}

\title{Analytical Estimate of Photon-Detector Backgrounds\\
ProtoDUNE-VD $\rightarrow$ DUNE FD-VD}
\author{}
\date{\today}

\begin{document}
\maketitle

\section{Purpose}

The objective is to estimate the photon-detector background seen by a cathode
X-ARAPUCA (XA) in ProtoDUNE-VD and extrapolate it to the DUNE vertical-drift
far detector.  The two dominant physical components considered here are
\begin{equation}
R_{\mathrm{tot}}
=
R_{39}
+
R_{\mu}
+
R_{\mathrm{other\,rad}}
+
R_{\mathrm{SiPM}},
\end{equation}
where $R_{39}$ is the contribution from $^{39}$Ar, $R_{\mu}$ is cosmic-muon
induced scintillation light, $R_{\mathrm{other\,rad}}$ contains other
radioactive sources, and $R_{\mathrm{SiPM}}$ includes detector dark noise and
correlated SiPM noise.

The useful comparison variable is
\begin{equation}
\boxed{R_{\mathrm{PE}}\quad [\mathrm{PE/s/XA}]},
\end{equation}
rather than total detector activity.

\section{Detector geometry used}

ProtoDUNE-VD has a central cathode with approximately \SI{3.5}{m} drift on
each side, with eight cathode X-ARAPUCAs.  The relevant transverse dimensions
are approximately
\begin{equation}
X_{\mathrm{PD}} \simeq \SI{6.8}{m},
\qquad
Y_{\mathrm{PD}} \simeq \SI{3.0}{m},
\qquad
L_{\mathrm{PD}} \simeq \SI{3.5}{m}.
\end{equation}

For DUNE FD-VD, the central cathode separates two drift regions with a drift
distance of approximately
\begin{equation}
L_{\mathrm{FD}} \simeq \SI{6.5}{m},
\end{equation}
and the cathode contains 320 double-sided X-ARAPUCAs.  A representative
transverse scale is
\begin{equation}
X_{\mathrm{FD}} \simeq \SI{13.5}{m},
\qquad
Y_{\mathrm{FD}} \simeq \SI{60}{m}.
\end{equation}

The nominal XA frontal area is
\begin{equation}
A_{\mathrm{XA}} \simeq
(0.60)(0.60)
=
\SI{0.36}{m^2}.
\end{equation}

\section{Geometrical acceptance for a uniform radioactive source}

Consider a point-like isotropic scintillation source at
$\bm{r}=(x,y,z)$ relative to the center of an XA.  In the small-area
approximation, the geometrical probability for a photon to intersect the XA is
\begin{equation}
P_{\mathrm{geom}}(x,y,z)
\simeq
\frac{A_{\mathrm{XA}}\cos\theta}{4\pi r^2}
=
\frac{A_{\mathrm{XA}}}{4\pi}
\frac{z}{(x^2+y^2+z^2)^{3/2}}.
\end{equation}

For a uniform source density, define
\begin{equation}
G
=
\int_V P_{\mathrm{geom}}\,dV.
\end{equation}

For one rectangular drift volume,
\begin{equation}
\frac{G_{\mathrm{side}}}{A_{\mathrm{XA}}}
=
\frac{1}{4\pi}
\int_0^L
\int_{-X/2}^{X/2}
\int_{-Y/2}^{Y/2}
\frac{z\,dx\,dy\,dz}
{(x^2+y^2+z^2)^{3/2}}.
\label{eq:Gvolume}
\end{equation}

The $x$--$y$ integral is the solid angle subtended by the rectangle:
\begin{equation}
\Omega(z)
=
4\tan^{-1}
\left[
\frac{ab}
{z\sqrt{z^2+a^2+b^2}}
\right],
\qquad
a=\frac{X}{2},
\quad
b=\frac{Y}{2}.
\end{equation}

Therefore,
\begin{equation}
\boxed{
\frac{G_{\mathrm{side}}}{A_{\mathrm{XA}}}
=
\frac{1}{4\pi}
\int_0^L \Omega(z)\,dz
}
\label{eq:Gside}
\end{equation}
and for a double-sided cathode XA,
\begin{equation}
\boxed{
\frac{G_{\mathrm{DS}}}{A_{\mathrm{XA}}}
=
2\frac{G_{\mathrm{side}}}{A_{\mathrm{XA}}}.
}
\end{equation}

In the infinite-plane limit, $\Omega(z)\rightarrow 2\pi$, giving
\begin{equation}
\frac{G_{\mathrm{DS}}}{A_{\mathrm{XA}}}
\rightarrow L.
\end{equation}

\section{Finite-geometry values}

Numerical integration of Eq.~\eqref{eq:Gside} gives approximately
\begin{align}
\left(\frac{G_{\mathrm{DS}}}{A_{\mathrm{XA}}}\right)_{\mathrm{PD}}
&\simeq
\SI{1.61}{m},
\\
\left(\frac{G_{\mathrm{DS}}}{A_{\mathrm{XA}}}\right)_{\mathrm{FD}}
&\simeq
\SI{4.69}{m}.
\end{align}

Thus, before optical-transport corrections,
\begin{equation}
\boxed{
\frac{R_{39}^{\mathrm{FD}}}
{R_{39}^{\mathrm{PD}}}
\simeq
\frac{4.69}{1.61}
\simeq
2.9.
}
\label{eq:ratio39}
\end{equation}

This is the key geometrical result: the $^{39}$Ar background per cathode XA
does \emph{not} scale with total detector mass.  It grows only by a factor of
order three between ProtoDUNE-VD and FD-VD in this first-order direct-light
model.

\section{$^{39}$Ar activity}

The measured atmospheric-argon specific activity is
\begin{equation}
S_{39}
=
(0.964 \pm 0.001_{\mathrm{stat}}
\pm 0.024_{\mathrm{sys}})
\ \mathrm{Bq/kg}.
\end{equation}

Using
\begin{equation}
\rho_{\mathrm{LAr}}\simeq \SI{1395}{kg/m^3},
\end{equation}
the volumetric decay density is
\begin{equation}
a_{39}
=
S_{39}\rho_{\mathrm{LAr}}
\simeq
\SI{1.34e3}{decays.s^{-1}.m^{-3}}.
\end{equation}

Define the effective geometrical volume
\begin{equation}
V_{\mathrm{eff}}
=
G_{\mathrm{DS}}
=
A_{\mathrm{XA}}
\left(
\frac{G_{\mathrm{DS}}}{A_{\mathrm{XA}}}
\right).
\end{equation}

With $A_{\mathrm{XA}}=\SI{0.36}{m^2}$,
\begin{align}
V_{\mathrm{eff}}^{\mathrm{PD}}
&\simeq
0.36\times1.61
\simeq
\SI{0.58}{m^3},
\\
V_{\mathrm{eff}}^{\mathrm{FD}}
&\simeq
0.36\times4.69
\simeq
\SI{1.69}{m^3}.
\end{align}

Hence the geometrically weighted $^{39}$Ar decay rates are
\begin{align}
R_{39,\mathrm{decay}}^{\mathrm{PD}}
&=
a_{39}V_{\mathrm{eff}}^{\mathrm{PD}}
\simeq
\SI{7.8e2}{s^{-1}},
\\
R_{39,\mathrm{decay}}^{\mathrm{FD}}
&=
a_{39}V_{\mathrm{eff}}^{\mathrm{FD}}
\simeq
\SI{2.3e3}{s^{-1}}.
\end{align}

If the mean scintillation yield per $^{39}$Ar decay is denoted by
$\overline{N}_{\gamma,39}$, then
\begin{equation}
R_{\gamma,39}
=
R_{39,\mathrm{decay}}
\,\overline{N}_{\gamma,39}.
\end{equation}

The corresponding detected PE rate is
\begin{equation}
\boxed{
R_{\mathrm{PE},39}
=
a_{39}
\,A_{\mathrm{XA}}
\left(
\frac{G_{\mathrm{DS}}}{A_{\mathrm{XA}}}
\right)
\overline{N}_{\gamma,39}
\,\epsilon_{\mathrm{eff}}.
}
\label{eq:R39}
\end{equation}

Here $\epsilon_{\mathrm{eff}}$ must be defined carefully.  If the quoted XA
efficiency already includes wavelength shifting, trapping, optical-window
transmission, SiPM PDE, and collection effects, those factors must not be
applied a second time.

For the working assumption
\begin{equation}
\epsilon_{\mathrm{eff}}\sim 0.05,
\end{equation}
Eq.~\eqref{eq:R39} provides the conversion from the geometrical estimate to
PE/s/XA once $\overline{N}_{\gamma,39}$ is fixed consistently with the drift
field and xenon concentration.

\section{Cosmic-muon light term}

For a single cosmic muon with trajectory $\Gamma$, the deposited energy is
\begin{equation}
E_{\mu}
=
\int_{\Gamma}
\left(\frac{dE}{dx}\right)_{\mu}
\,d\ell.
\end{equation}

The emitted scintillation light is
\begin{equation}
N_{\gamma,\mu}
=
Y_{\gamma}(E,\mathcal{E},c_{\mathrm{Xe}})
\,E_{\mu},
\end{equation}
where $Y_{\gamma}$ depends on deposited-energy density, electric field
$\mathcal{E}$, and xenon concentration $c_{\mathrm{Xe}}$.

For XA $j$, introduce the optical transfer function
\begin{equation}
T_j(\bm{r})
=
P_{\mathrm{geom},j}(\bm{r})
\,
P_{\mathrm{transport}}(\bm{r}\rightarrow j)
\,
\epsilon_{\mathrm{XA},j}.
\end{equation}

Then the PE yield for one muon is
\begin{equation}
N_{\mathrm{PE},\mu}^{(j)}
=
\int_{\Gamma}
Y_{\gamma}
\left(\frac{dE}{dx}\right)_{\mu}
T_j(\bm{r})
\,d\ell.
\label{eq:muPE}
\end{equation}

For a cosmic-muon differential flux
$\Phi_{\mu}(\theta,\phi)$, the time-averaged rate is
\begin{equation}
\boxed{
R_{\mathrm{PE},\mu}^{(j)}
=
\int
d\Omega
\int
dA_{\perp}\,
\Phi_{\mu}(\theta,\phi)
\,
N_{\mathrm{PE},\mu}^{(j)}
(\theta,\phi,\bm{r}_0).
}
\label{eq:muRate}
\end{equation}

Experimentally, the cleaner estimator is obtained directly from CRT-tagged
events:
\begin{equation}
\boxed{
R_{\mathrm{PE},\mu}^{(j)}
=
f_{\mu,\mathrm{CRT}}
\,
\left\langle
N_{\mathrm{PE},\mu}^{(j)}
\right\rangle,
}
\label{eq:CRT}
\end{equation}
after correcting for CRT geometrical acceptance and efficiency.

\section{ProtoDUNE-VD to FD-VD extrapolation}

Define
\begin{equation}
R_{\mathrm{PD}}
=
R_{39}^{\mathrm{PD}}
+
R_{\mu}^{\mathrm{PD}}
+
R_{\mathrm{other}}^{\mathrm{PD}},
\end{equation}
and
\begin{equation}
R_{\mathrm{FD}}
=
R_{39}^{\mathrm{FD}}
+
R_{\mu}^{\mathrm{FD}}
+
R_{\mathrm{other}}^{\mathrm{FD}}.
\end{equation}

The first-order radiological extrapolation is
\begin{equation}
\boxed{
R_{39}^{\mathrm{FD}}
\simeq
2.9\,
R_{39}^{\mathrm{PD}},
}
\end{equation}
subject to changes in optical transport, Xe doping, electric field,
reflectivity, detector response, and active optical area.

The cosmic term should instead scale approximately as
\begin{equation}
R_{\mu}^{\mathrm{FD}}
=
R_{\mu}^{\mathrm{PD}}
\,
\frac{\Phi_{\mu}^{\mathrm{FD}}}
{\Phi_{\mu}^{\mathrm{PD}}}
\,
C_{\mathrm{geom}}
\,
C_{\mathrm{opt}},
\end{equation}
where $C_{\mathrm{geom}}$ accounts for the different detector dimensions and
track-length distributions, and $C_{\mathrm{opt}}$ accounts for different
light transport.

The core expectation is therefore
\begin{equation}
\boxed{
\text{ProtoDUNE-VD: cosmic bursts + radiological floor}
\quad\longrightarrow\quad
\text{FD-VD: radiological floor dominates}.
}
\end{equation}

\section{Quantities to extract from ProtoDUNE-VD data}

The most useful experimental observables are
\begin{enumerate}
    \item random-trigger PE-rate distribution per cathode XA;
    \item CRT-tagged PE per muon per cathode XA;
    \item muon position and direction relative to each XA;
    \item off-CRT background in the same acquisition windows;
    \item dependence on Xe concentration;
    \item SiPM dark-noise and correlated-noise contribution;
    \item channel-to-channel variation in effective PDE.
\end{enumerate}

A direct decomposition can then be written as
\begin{equation}
R_{\mathrm{off}}
=
R_{39}
+
R_{\mathrm{other\,rad}}
+
R_{\mathrm{SiPM}},
\end{equation}
and
\begin{equation}
R_{\mathrm{on}}-R_{\mathrm{off}}
\rightarrow
R_{\mu}.
\end{equation}

This allows the measured ProtoDUNE-VD background to be extrapolated to FD-VD
without relying exclusively on full optical Monte Carlo.

\section{Important limitations}

The present model is intentionally analytical.  It neglects or compresses
into effective factors:
\begin{itemize}
    \item Rayleigh scattering;
    \item VUV absorption;
    \item reflections from detector structures;
    \item field-cage transparency;
    \item wavelength shifting and Xe transfer;
    \item angular dependence of XA response;
    \item finite XA active-window fraction;
    \item nonuniform detector response;
    \item other radiological isotopes and detector materials.
\end{itemize}

Therefore, Eq.~\eqref{eq:ratio39} should be interpreted as a
\emph{finite-geometry direct-light scaling estimate}, not yet as a final
prediction of measured PE rate.

\section{References}

\begin{thebibliography}{9}

\bibitem{pdvd_geometry}
L. Zambelli,
\emph{ProtoDUNE-VD: Module-0},
NuFact 2023 presentation.
\url{https://indico.cern.ch/event/1216905/contributions/5456597/attachments/2700868/4688179/lzambell_nufact23_protodunevd.pdf}

\bibitem{fdvd_pde}
\emph{Measurement of the X-ARAPUCA's absolute photon detection efficiency for
the Deep Underground Neutrino Experiment's vertical drift far detector},
Eur. Phys. J. C (2026).
\url{https://link.springer.com/article/10.1140/epjc/s10052-026-16190-1}

\bibitem{ar39}
DEAP-3600 Collaboration,
\emph{Precision measurement of the specific activity of $^{39}$Ar in
atmospheric argon with the DEAP-3600 detector},
Eur. Phys. J. C 83 (2023).
\url{https://link.springer.com/article/10.1140/epjc/s10052-023-11678-6}

\end{thebibliography}

\end{document}
